% ============================================================
% CHAPITRE 18 — Couleurs, dégradés et contours
% ============================================================

\chapter{Couleurs, dégradés et contours}

\noindent
\textit{Dans les chapitres précédents, nous avons appris à dessiner avec Illustrator : formes géométriques, plume, courbes de Bézier, manipulation avancée des formes. Nous savons désormais créer des tracés précis, combiner des formes, utiliser les masques, organiser les objets. Mais un dessin sans couleur reste un dessin en noir et blanc. C'est la couleur, le dégradé, le contour qui donnent vie à une illustration, qui transmettent une émotion, qui créent une identité.}

\vspace{1em}

\noindent
La \terme{couleur} est au cœur du design graphique. Dans Illustrator, elle s'applique aux objets de deux manières : le \terme{remplissage} (l'intérieur de la forme) et le \terme{contour} (la ligne qui dessine la forme). On peut utiliser des couleurs unies, des dégradés, des motifs, des transparences. On peut enregistrer des couleurs globales, créer des palettes cohérentes, préparer les couleurs pour l'impression. Maîtriser la couleur dans Illustrator, c'est maîtriser l'un des leviers les plus puissants du design graphique.

\vspace{1em}

\noindent
Mais la couleur ne s'utilise pas au hasard. Elle doit respecter les principes étudiés dans la première partie de ce cours (théorie des couleurs, harmonies, psychologie, contraste, lisibilité). Elle doit être adaptée au message, au public, au support. Elle doit être cohérente sur l'ensemble d'un projet. Elle doit être préparée correctement pour l'impression ou pour l'écran. C'est cette dimension professionnelle que nous allons explorer dans ce chapitre.

\vspace{1em}

\noindent
Ce chapitre a pour objectif de te faire maîtriser les couleurs, les dégradés et les contours dans Illustrator. Nous commencerons par apprendre à appliquer des couleurs aux objets. Nous étudierons ensuite les nuanciers, la création et l'enregistrement de couleurs globales, les dégradés linéaires et radiaux, la maîtrise des contours, le réglage des extrémités et des angles, la création de motifs répétitifs, la gestion des transparences, et enfin la préparation des couleurs pour l'impression. À la fin de ce chapitre, tu sauras utiliser la couleur de manière professionnelle.

% ============================================================
\section{Appliquer des couleurs aux objets}
% ============================================================

\subsection{Le remplissage et le contour}

\subsubsection{Les deux propriétés}

Dans Illustrator, chaque objet possède deux propriétés de couleur :

\begin{itemize}
    \item \terme{Le remplissage} (ou \emph{fill}) : la couleur, le dégradé ou le motif qui remplit l'intérieur de la forme.
    \item \terme{Le contour} (ou \emph{stroke}) : la couleur, le dégradé ou le motif de la ligne qui dessine la forme.
\end{itemize}

Chaque propriété peut être :

\begin{itemize}
    \item \terme{Aucune} : pas de remplissage (transparent) ou pas de contour.
    \item \terme{Couleur unie} : une couleur unique.
    \item \terme{Dégradé} : une transition entre plusieurs couleurs.
    \item \terme{Motif} : un motif répété.
\end{itemize}

\subsubsection{La case Remplissage et la case Contour}

Dans la barre d'outils, deux cases superposées représentent le remplissage et le contour :

\begin{itemize}
    \item \terme{La case du haut} : le remplissage.
    \item \terme{La case du bas} : le contour.
    \item \terme{L'icône de permutation} (flèche courbe) : permute le remplissage et le contour (\texttt{Maj + X}).
    \item \terme{L'icône par défaut} (petit carré noir/blanc) : remet le remplissage en blanc et le contour en noir (\texttt{D}).
    \item \terme{L'icône Aucun} (carré rouge barré) : supprime le remplissage ou le contour (\texttt{/}).
\end{itemize}

\subsection{Appliquer une couleur unie}

\subsubsection{Avec le panneau Couleur}

\begin{enumerate}
    \item Sélectionner l'objet.
    \item Ouvrir le panneau \emph{Couleur} (\texttt{Fenêtre > Couleur}).
    \item Choisir la couleur (RVB, CMJN, TSL, etc.).
    \item Le remplissage (ou le contour, selon la case active) est appliqué.
\end{enumerate}

\subsubsection{Avec le panneau Nuancier}

\begin{enumerate}
    \item Sélectionner l'objet.
    \item Ouvrir le panneau \emph{Nuancier} (\texttt{Fenêtre > Nuancier}).
    \item Cliquer sur une couleur.
    \item Le remplissage (ou le contour) est appliqué.
\end{enumerate}

\subsubsection{Avec le sélecteur de couleur}

\begin{enumerate}
    \item Double-cliquer sur la case de remplissage ou de contour dans la barre d'outils.
    \item Le sélecteur de couleur s'ouvre.
    \item Choisir une couleur.
    \item Valider.
\end{enumerate}

\subsection{Appliquer un contour}

\subsubsection{Les propriétés du contour}

Le contour a plusieurs propriétés :

\begin{itemize}
    \item \terme{Couleur} : la couleur du contour.
    \item \terme{Épaisseur} : l'épaisseur du contour (en points).
    \item \terme{Style} : le style du contour (continu, pointillé, tiret).
    \item \terme{Extrémité} : le bout du contour (bout, rond, carré).
    \item \terme{Angle} : le coin du contour (angle, arrondi, biseau).
\end{itemize}

\subsubsection{Appliquer une épaisseur de contour}

\begin{enumerate}
    \item Sélectionner l'objet.
    \item Ouvrir le panneau \emph{Contour} (\texttt{Fenêtre > Contour}).
    \item Saisir une épaisseur (en points).
    \item Le contour est appliqué.
\end{enumerate}

\subsection{Les raccourcis}

\begin{itemize}
    \item \terme{X} : permuter le remplissage et le contour.
    \item \terme{D} : couleurs par défaut (blanc/noir).
    \item \terme{/} : aucun remplissage ou contour.
    \item \terme{Maj + X} : permuter les couleurs.
    \item \terme{<} : appliquer la couleur du nuancier au remplissage (si sélectionné).
    \item \terme{>} : appliquer la couleur du nuancier au contour.
\end{itemize}

\subsection{Les bonnes pratiques}

\begin{itemize}
    \item \terme{Comprendre} la différence entre remplissage et contour.
    \item \terme{Utiliser} les raccourcis (\texttt{X}, \texttt{D}, \texttt{/}).
    \item \terme{Choisir} les couleurs dans le nuancier.
    \item \terme{Vérifier} le résultat.
    \item \terme{Travailler} avec des couleurs cohérentes.
\end{itemize}

\subsection*{Conclusion de la section}

Chaque objet dans Illustrator possède un \terme{remplissage} et un \terme{contour}. On peut leur appliquer des couleurs unies, des dégradés, des motifs, ou aucun. On utilise le panneau Couleur, le panneau Nuancier, ou le sélecteur de couleur. Les raccourcis (\texttt{X}, \texttt{D}, \texttt{/}) permettent de manipuler rapidement les couleurs. Il est essentiel de comprendre la différence entre remplissage et contour, d'utiliser les raccourcis, de choisir des couleurs cohérentes.

Maintenant que nous savons appliquer des couleurs, nous allons, dans la section suivante, étudier les \terme{nuanciers}.

% ============================================================
\section{Utiliser les nuanciers}
% ============================================================

\subsection{Qu'est-ce qu'un nuancier ?}

\begin{quote}
\textbf{Un nuancier est une bibliothèque de couleurs prédéfinies ou enregistrées, organisée pour être facilement accessible dans Illustrator.}
\end{quote}

Le nuancier est un outil essentiel pour :

\begin{itemize}
    \item \terme{Accéder} rapidement aux couleurs.
    \item \terme{Maintenir} la cohérence des couleurs.
    \item \terme{Organiser} les couleurs d'un projet.
    \item \terme{Partager} des couleurs avec d'autres designers.
\end{itemize}

\subsection{Le panneau Nuancier}

\subsubsection{Accès au panneau}

Le panneau \terme{Nuancier} (\texttt{Fenêtre > Nuancier}) affiche une grille de couleurs.

\subsubsection{Les types de nuanciers}

\begin{itemize}
    \item \terme{Couleurs unies} : couleurs uniques.
    \item \terme{Dégradés} : dégradés enregistrés.
    \item \terme{Motifs} : motifs enregistrés.
    \item \terme{Groupes de couleurs} : groupes de couleurs.
\end{itemize}

\subsubsection{Les bibliothèques}

Illustrator propose des \terme{bibliothèques} de nuanciers :

\begin{itemize}
    \item \terme{Default Swatches} : couleurs par défaut.
    \item \terme{Pantone} : couleurs Pantone (pour l'impression).
    \item \terme{Web} : couleurs pour le web.
    \item \terme{Print} : couleurs pour l'impression.
    \item \terme{Thèmes} : palettes thématiques.
\end{itemize}

On peut ouvrir une bibliothèque via \texttt{Fenêtre > Bibliothèques de nuanciers}.

\subsection{Utiliser le nuancier}

\subsubsection{Appliquer une couleur}

\begin{enumerate}
    \item Sélectionner l'objet.
    \item Sélectionner la case de remplissage ou de contour (selon la cible).
    \item Cliquer sur une couleur dans le nuancier.
    \item La couleur est appliquée.
\end{enumerate}

\subsubsection{Créer une couleur}

\begin{enumerate}
    \item Cliquer sur \emph{Nouvelle couleur} dans le panneau Nuancier.
    \item Choisir la couleur.
    \item Nommer la couleur.
    \item Valider.
\end{enumerate}

\subsubsection{Supprimer une couleur}

\begin{enumerate}
    \item Sélectionner la couleur dans le nuancier.
    \item Cliquer sur \emph{Supprimer} (icône de corbeille).
\end{enumerate}

\subsection{Les groupes de couleurs}

\subsubsection{Créer un groupe}

\begin{enumerate}
    \item Sélectionner plusieurs couleurs (avec \texttt{Ctrl} ou \texttt{Maj}).
    \item Cliquer sur \emph{Nouveau groupe de couleurs}.
    \item Nommer le groupe.
    \item Valider.
\end{enumerate}

\subsubsection{Utilisation}

Les groupes de couleurs permettent de :

\begin{itemize}
    \item \terme{Organiser} les couleurs d'un projet.
    \item \terme{Modifier} plusieurs couleurs en même temps.
    \item \terme{Partager} des palettes.
\end{itemize}

\subsection{Les bonnes pratiques}

\begin{itemize}
    \item \terme{Utiliser} les nuanciers pour la cohérence.
    \item \terme{Créer} ses propres couleurs.
    \item \terme{Organiser} les couleurs en groupes.
    \item \terme{Enregistrer} les palettes.
    \item \terme{Utiliser} les bibliothèques prédéfinies.
    \item \terme{Vérifier} les couleurs avant l'export.
\end{itemize}

\subsection*{Conclusion de la section}

Un \terme{nuancier} est une bibliothèque de couleurs. On l'utilise pour accéder rapidement aux couleurs, maintenir la cohérence, organiser les couleurs d'un projet. On peut créer, supprimer, organiser des couleurs, et utiliser des bibliothèques prédéfinies (Pantone, Web, Print). Il est essentiel d'utiliser les nuanciers, de créer ses propres couleurs, d'organiser en groupes.

Maintenant que nous connaissons les nuanciers, nous allons, dans la section suivante, apprendre à \terme{créer et enregistrer des couleurs globales}.

% ============================================================
\section{Créer et enregistrer des couleurs globales}
% ============================================================

\subsection{Qu'est-ce qu'une couleur globale ?}

\begin{quote}
\textbf{Une couleur globale (ou \emph{global color}) est une couleur enregistrée dans le nuancier, qui peut être modifiée et mise à jour automatiquement dans tous les objets qui l'utilisent.}
\end{quote}

Les couleurs globales sont essentielles pour :

\begin{itemize}
    \item \terme{Maintenir} la cohérence des couleurs.
    \item \terme{Modifier} toutes les occurrences d'une couleur en une seule fois.
    \item \terme{Préparer} les couleurs pour l'impression.
    \item \terme{Travailler} avec des palettes évolutives.
\end{itemize}

\subsection{Créer une couleur globale}

\subsubsection{Par le nuancier}

\begin{enumerate}
    \item Ouvrir le panneau \emph{Nuancier}.
    \item Cliquer sur \emph{Nouvelle couleur globale}.
    \item Choisir la couleur.
    \item Nommer la couleur.
    \item Valider.
\end{enumerate}

\subsubsection{Par le sélecteur de couleur}

\begin{enumerate}
    \item Double-cliquer sur la case de remplissage ou de contour.
    \item Choisir une couleur.
    \item Cocher \emph{Couleur globale}.
    \item Valider.
    \item La couleur est ajoutée au nuancier.
\end{enumerate}

\subsubsection{Convertir une couleur en couleur globale}

\begin{enumerate}
    \item Double-cliquer sur une couleur dans le nuancier.
    \item Cocher \emph{Couleur globale}.
    \item Valider.
\end{enumerate}

\subsection{Modifier une couleur globale}

\begin{enumerate}
    \item Double-cliquer sur la couleur globale dans le nuancier.
    \item Modifier la couleur.
    \item Valider.
    \item Tous les objets utilisant cette couleur sont automatiquement mis à jour.
\end{enumerate}

\subsection{Les avantages}

\begin{itemize}
    \item \terme{Cohérence} : toutes les occurrences sont identiques.
    \item \terme{Modification} : un seul changement met à jour tous les objets.
    \item \terme{Organisation} : les couleurs sont nommées et regroupées.
    \item \terme{Impression} : les couleurs peuvent être converties en tons directs.
\end{itemize}

\subsection{Les couleurs de tons directs}

\subsubsection{Définition}

Un \terme{ton direct} (ou \emph{spot color}) est une couleur d'encre spéciale, imprimée avec une plaque dédiée, en plus des quatre encres CMJN. Les couleurs Pantone sont des tons directs.

\subsubsection{Utilisation}

Les tons directs sont utilisés pour :

\begin{itemize}
    \item \terme{Les logos} : couleurs précises et constantes.
    \item \terme{Les identités visuelles} : couleurs de marque.
    \item \terme{Les impressions} : couleurs métalliques, fluo.
\end{itemize}

\subsection{Les bonnes pratiques}

\begin{itemize}
    \item \terme{Utiliser} les couleurs globales pour la cohérence.
    \item \terme{Nommer} les couleurs globales.
    \item \terme{Organiser} les couleurs en groupes.
    \item \terme{Vérifier} les couleurs avant l'export.
    \item \terme{Convertir} les couleurs de marque en tons directs.
    \item \terme{Enregistrer} les palettes.
\end{itemize}

\subsection*{Conclusion de la section}

Les \terme{couleurs globales} sont des couleurs enregistrées dans le nuancier, qui peuvent être modifiées et mises à jour automatiquement dans tous les objets. On les crée avec le nuancier ou le sélecteur de couleur. On peut les modifier, les convertir en tons directs. Les couleurs globales permettent de maintenir la cohérence, de modifier facilement, d'organiser, de préparer pour l'impression.

Maintenant que nous savons créer des couleurs globales, nous allons, dans la section suivante, étudier les \terme{dégradés linéaires}.

% ============================================================
\section{Utiliser les dégradés linéaires}
% ============================================================

\subsection{Qu'est-ce qu'un dégradé ?}

\begin{quote}
\textbf{Un dégradé est une transition progressive entre plusieurs couleurs, dans une direction donnée. Il permet de créer des effets de lumière, de profondeur, de volume.}
\end{quote}

Les dégradés sont utilisés pour :

\begin{itemize}
    \item \terme{Créer} des effets de lumière.
    \item \terme{Donner} du volume aux formes.
    \item \terme{Créer} des fonds.
    \item \terme{Ajouter} de la profondeur.
    \item \terme{Enrichir} les illustrations.
\end{itemize}

\subsection{Le panneau Dégradé}

\subsubsection{Accès au panneau}

Le panneau \terme{Dégradé} (\texttt{Fenêtre > Dégradé}) permet de créer et de modifier les dégradés.

\subsubsection{Les types de dégradés}

\begin{itemize}
    \item \terme{Linéaire} : transition dans une direction droite.
    \item \terme{Radial} : transition depuis un point central.
\end{itemize}

\subsection{Créer un dégradé linéaire}

\subsubsection{Avec le panneau Dégradé}

\begin{enumerate}
    \item Sélectionner l'objet.
    \item Ouvrir le panneau \emph{Dégradé}.
    \item Choisir le type \emph{Linéaire}.
    \item Cliquer sur la barre de dégradé pour ajouter des arrêts de couleur.
    \item Choisir les couleurs de chaque arrêt.
    \item Ajuster la position des arrêts.
    \item Ajuster l'angle du dégradé.
\end{enumerate}

\subsubsection{Avec l'outil Dégradé}

\begin{enumerate}
    \item Sélectionner l'objet.
    \item Sélectionner l'outil Dégradé (\texttt{G}).
    \item Cliquer-glisser sur l'objet pour définir la direction et la longueur du dégradé.
    \item Le dégradé est appliqué.
\end{enumerate}

\subsection{Modifier un dégradé}

\subsubsection{Ajouter un arrêt}

\begin{enumerate}
    \item Cliquer sur la barre de dégradé (dans le panneau Dégradé).
    \item Un nouvel arrêt est ajouté.
    \item Choisir la couleur de l'arrêt.
\end{enumerate}

\subsubsection{Déplacer un arrêt}

\begin{enumerate}
    \item Cliquer sur un arrêt.
    \item Glisser l'arrêt sur la barre.
    \item La transition est modifiée.
\end{enumerate}

\subsubsection{Supprimer un arrêt}

\begin{enumerate}
    \item Cliquer sur un arrêt.
    \item Glisser l'arrêt en dehors de la barre.
    \item L'arrêt est supprimé.
\end{enumerate}

\subsubsection{Ajuster l'angle}

\begin{enumerate}
    \item Dans le panneau Dégradé, saisir un angle.
    \item Le dégradé est orienté.
\end{enumerate}

\subsection{Les dégradés avec transparence}

\subsubsection{Le principe}

Un dégradé peut inclure de la transparence : un arrêt peut avoir une opacité inférieure à 100\%.

\subsubsection{Utilisation}

\begin{enumerate}
    \item Cliquer sur un arrêt.
    \item Changer l'opacité dans le panneau Dégradé.
    \item Le dégradé inclut la transparence.
\end{enumerate}

\subsection{Les bonnes pratiques}

\begin{itemize}
    \item \terme{Choisir} des couleurs harmonieuses.
    \item \terme{Limiter} le nombre d'arrêts.
    \item \terme{Ajuster} la position des arrêts.
    \item \terme{Utiliser} l'outil Dégradé pour des dégradés précis.
    \item \terme{Tester} le résultat.
    \item \terme{Enregistrer} les dégradés dans le nuancier.
\end{itemize}

\subsection*{Conclusion de la section}

Un \terme{dégradé} est une transition progressive entre plusieurs couleurs. Le dégradé \terme{linéaire} se fait dans une direction droite. On le crée avec le panneau Dégradé ou l'outil Dégradé (\texttt{G}). On peut ajouter, déplacer, supprimer des arrêts, ajuster l'angle, inclure de la transparence. Il est essentiel de choisir des couleurs harmonieuses, de limiter le nombre d'arrêts, de tester le résultat.

Maintenant que nous savons créer des dégradés linéaires, nous allons, dans la section suivante, étudier les \terme{dégradés radiaux}.

% ============================================================
\section{Utiliser les dégradés radiaux}
% ============================================================

\subsection{Qu'est-ce qu'un dégradé radial ?}

\begin{quote}
\textbf{Un dégradé radial est une transition progressive entre plusieurs couleurs, depuis un point central vers l'extérieur. Il permet de créer des effets de sphère, de halo, de lumière.}
\end{quote}

Les dégradés radiaux sont utilisés pour :

\begin{itemize}
    \item \terme{Créer} des sphères.
    \item \terme{Donner} du volume aux formes.
    \item \terme{Créer} des halos lumineux.
    \item \terme{Ajouter} de la profondeur.
    \item \terme{Créer} des fonds.
\end{itemize}

\subsection{Créer un dégradé radial}

\subsubsection{Avec le panneau Dégradé}

\begin{enumerate}
    \item Sélectionner l'objet.
    \item Ouvrir le panneau \emph{Dégradé}.
    \item Choisir le type \emph{Radial}.
    \item Cliquer sur la barre de dégradé pour ajouter des arrêts.
    \item Choisir les couleurs de chaque arrêt.
    \item Ajuster la position des arrêts.
\end{enumerate}

\subsubsection{Avec l'outil Dégradé}

\begin{enumerate}
    \item Sélectionner l'objet.
    \item Sélectionner l'outil Dégradé (\texttt{G}).
    \item Cliquer-glisser sur l'objet pour définir le centre et le rayon du dégradé.
    \item Le dégradé est appliqué.
\end{enumerate}

\subsection{Modifier un dégradé radial}

\subsubsection{Le centre}

Le centre du dégradé radial est le point où commence la transition. On peut le déplacer avec l'outil Dégradé.

\subsubsection{Le rayon}

Le rayon du dégradé radial est la distance sur laquelle la transition s'effectue. On peut l'ajuster avec l'outil Dégradé.

\subsubsection{L'aspect}

L'aspect du dégradé radial peut être modifié pour créer une ellipse au lieu d'un cercle.

\subsection{Les utilisations}

Les dégradés radiaux sont utilisés pour :

\begin{itemize}
    \item \terme{Sphères} : créer une impression de volume.
    \item \terme{Halos} : créer une lumière autour d'un objet.
    \item \terme{Fonds} : créer une ambiance.
    \item \terme{Yeux} : créer un effet de profondeur.
    \item \terme{Boutons} : créer un effet 3D.
\end{itemize}

\subsection{Les bonnes pratiques}

\begin{itemize}
    \item \terme{Choisir} des couleurs harmonieuses.
    \item \terme{Ajuster} le centre et le rayon.
    \item \terme{Limiter} le nombre d'arrêts.
    \item \terme{Utiliser} l'outil Dégradé pour des dégradés précis.
    \item \terme{Tester} le résultat.
    \item \terme{Enregistrer} les dégradés dans le nuancier.
\end{itemize}

\subsection*{Conclusion de la section}

Un \terme{dégradé radial} est une transition progressive depuis un point central vers l'extérieur. On le crée avec le panneau Dégradé ou l'outil Dégradé (\texttt{G}). On peut ajuster le centre, le rayon, l'aspect. Les dégradés radiaux sont utilisés pour créer des sphères, des halos, des fonds, des effets de profondeur. Il est essentiel de choisir des couleurs harmonieuses, d'ajuster le centre et le rayon, de tester le résultat.

Maintenant que nous savons créer des dégradés radiaux, nous allons, dans la section suivante, apprendre à \terme{maîtriser les contours}.

% ============================================================
\section{Maîtriser les contours}
% ============================================================

\subsection{Le rôle du contour}

Le \terme{contour} (ou \emph{stroke}) est la ligne visible qui dessine un tracé. Il a plusieurs rôles :

\begin{itemize}
    \item \terme{Délimiter} une forme.
    \item \terme{Détacher} une forme du fond.
    \item \terme{Structurer} une composition.
    \item \terme{Décorer} une forme.
    \item \terme{Signifier} (par exemple, un contour épais pour un message fort).
\end{itemize}

\subsection{Le panneau Contour}

\subsubsection{Accès au panneau}

Le panneau \terme{Contour} (\texttt{Fenêtre > Contour}) permet de modifier les propriétés du contour.

\subsubsection{Les propriétés}

\begin{itemize}
    \item \terme{Épaisseur} : l'épaisseur du contour (en points).
    \item \terme{Extrémité} : le bout du contour (bout, rond, carré).
    \item \terme{Angle} : le coin du contour (angle, arrondi, biseau).
    \item \terme{Profil} : le profil du contour (uniforme, triangulaire, etc.).
    \item \terme{Pointillés} : pointillés (dash), espaces (gap).
    \item \terme{Flèches} : flèches de début et de fin.
\end{itemize}

\subsection{Modifier l'épaisseur}

\subsubsection{Épaisseur uniforme}

\begin{enumerate}
    \item Sélectionner l'objet.
    \item Ouvrir le panneau \emph{Contour}.
    \item Saisir une épaisseur (en points).
    \item Le contour est appliqué.
\end{enumerate}

\subsubsection{Épaisseur variable}

\begin{enumerate}
    \item Sélectionner l'objet.
    \item Utiliser l'outil \emph{Largeur} (dans le groupe de la Plume).
    \item Cliquer-glisser sur le contour pour modifier l'épaisseur localement.
    \item Le contour a une épaisseur variable.
\end{enumerate}

\subsection{Les pointillés}

\subsubsection{Créer des pointillés}

\begin{enumerate}
    \item Sélectionner l'objet.
    \item Ouvrir le panneau \emph{Contour}.
    \item Cocher \emph{Pointillés}.
    \item Saisir la longueur des tirets et l'espacement.
    \item Le contour est en pointillés.
\end{enumerate}

\subsubsection{Les styles de pointillés}

\begin{itemize}
    \item \terme{Tirets} : traits courts et espaces.
    \item \terme{Points} : points et espaces.
    \item \terme{Tirets-points} : alternance de tirets et de points.
    \item \terme{Personnalisés} : motifs personnalisés.
\end{itemize}

\subsection{Le profil de contour}

\subsubsection{Définition}

Le \terme{profil de contour} définit la variation d'épaisseur le long d'un tracé. Il permet de créer des contours dont l'épaisseur change progressivement.

\subsubsection{Les profils prédéfinis}

\begin{itemize}
    \item \terme{Uniforme} : épaisseur constante.
    \item \terme{Triangulaire} : épaisseur qui augmente puis diminue.
    \item \terme{Fuselé} : épaisseur qui diminue aux extrémités.
    \item \terme{Personnalisé} : profil personnalisé.
\end{itemize}

\subsection{Les bonnes pratiques}

\begin{itemize}
    \item \terme{Choisir} l'épaisseur adaptée.
    \item \terme{Utiliser} l'outil Largeur pour les contours variables.
    \item \terme{Créer} des pointillés pour des effets.
    \item \terme{Utiliser} les profils pour des contours expressifs.
    \item \terme{Vérifier} le résultat.
    \item \terme{Enregistrer} les styles dans les styles graphiques.
\end{itemize}

\subsection*{Conclusion de la section}

Le \terme{contour} est la ligne visible qui dessine un tracé. Il a plusieurs propriétés : épaisseur, extrémité, angle, profil, pointillés, flèches. On le modifie via le panneau Contour. On peut créer des contours d'épaisseur uniforme ou variable, des pointillés, des profils personnalisés. Il est essentiel de choisir l'épaisseur adaptée, d'utiliser l'outil Largeur, de créer des pointillés, de vérifier le résultat.

Maintenant que nous savons maîtriser les contours, nous allons, dans la section suivante, apprendre à \terme{régler les extrémités et les angles des contours}.

% ============================================================
\section{Régler les extrémités et les angles des contours}
% ============================================================

\subsection{Les extrémités de contour}

\subsubsection{Définition}

Les \terme{extrémités} (ou \emph{caps}) définissent l'apparence du bout d'un contour ouvert. Elles peuvent être :

\begin{itemize}
    \item \terme{Bout} (\emph{butt}) : le contour se termine net, sans extension.
    \item \terme{Rond} (\emph{round}) : le contour se termine par un demi-cercle.
    \item \terme{Carré} (\emph{square}) : le contour se termine par un carré.
\end{itemize}

\subsubsection{Utilisation}

\begin{enumerate}
    \item Sélectionner l'objet.
    \item Ouvrir le panneau \emph{Contour}.
    \item Choisir l'extrémité dans le menu déroulant.
    \item Le contour est modifié.
\end{enumerate}

\subsection{Les angles de contour}

\subsubsection{Définition}

Les \terme{angles} (ou \emph{joins}) définissent l'apparence des coins d'un contour. Ils peuvent être :

\begin{itemize}
    \item \terme{Angle} (\emph{miter}) : le coin forme une pointe.
    \item \terme{Arrondi} (\emph{round}) : le coin est arrondi.
    \item \terme{Biseau} (\emph{bevel}) : le coin est coupé.
\end{itemize}

\subsubsection{Utilisation}

\begin{enumerate}
    \item Sélectionner l'objet.
    \item Ouvrir le panneau \emph{Contour}.
    \item Choisir l'angle dans le menu déroulant.
    \item Le contour est modifié.
\end{enumerate}

\subsection{Le paramètre de limite d'angle}

\subsubsection{Définition}

Le \terme{paramètre de limite d'angle} (ou \emph{miter limit}) définit la longueur maximale d'une pointe d'angle. Si la pointe est trop longue, elle est remplacée par un biseau.

\subsubsection{Utilisation}

\begin{enumerate}
    \item Sélectionner l'objet.
    \item Ouvrir le panneau \emph{Contour}.
    \item Saisir une valeur dans le champ \emph{Limite d'angle}.
    \item Le contour est modifié.
\end{enumerate}

\subsection{Les flèches}

\subsubsection{Définition}

Les \terme{flèches} sont des symboles placés aux extrémités d'un contour ouvert. Elles peuvent être :

\begin{itemize}
    \item \terme{Flèches simples} : triangles.
    \item \terme{Flèches doubles} : triangles inversés.
    \item \terme{Points} : cercles.
    \item \terme{Barres} : traits.
    \item \terme{Personnalisées} : symboles personnalisés.
\end{itemize}

\subsubsection{Utilisation}

\begin{enumerate}
    \item Sélectionner l'objet.
    \item Ouvrir le panneau \emph{Contour}.
    \item Choisir la flèche de début et/ou de fin.
    \item Ajuster l'échelle.
    \item Le contour est modifié.
\end{enumerate}

\subsection{Les bonnes pratiques}

\begin{itemize}
    \item \terme{Choisir} l'extrémité adaptée.
    \item \terme{Choisir} l'angle adapté.
    \item \terme{Ajuster} la limite d'angle si nécessaire.
    \item \terme{Utiliser} les flèches pour les tracés directionnels.
    \item \terme{Vérifier} le résultat.
    \item \terme{Enregistrer} les styles dans les styles graphiques.
\end{itemize}

\subsection*{Conclusion de la section}

Les \terme{extrémités} (bout, rond, carré) et les \terme{angles} (angle, arrondi, biseau) définissent l'apparence des contours. On peut ajuster la limite d'angle, ajouter des flèches. On les modifie via le panneau Contour. Il est essentiel de choisir l'extrémité et l'angle adaptés, d'ajuster la limite d'angle, de vérifier le résultat.

Maintenant que nous savons régler les extrémités et les angles, nous allons, dans la section suivante, apprendre à \terme{créer des motifs répétitifs}.

% ============================================================
\section{Créer des motifs répétitifs}
% ============================================================

\subsection{Qu'est-ce qu'un motif ?}

\begin{quote}
\textbf{Un motif (ou \emph{pattern}) est une image répétée pour créer un fond, une texture, une décoration. Il peut être appliqué comme remplissage ou comme contour.}
\end{quote}

Les motifs sont utilisés pour :

\begin{itemize}
    \item \terme{Créer} des fonds.
    \item \terme{Ajouter} des textures.
    \item \terme{Décorer} des formes.
    \item \terme{Enrichir} les illustrations.
    \item \terme{Créer} des identités visuelles.
\end{itemize}

\subsection{Créer un motif}

\subsubsection{Méthode de base}

\begin{enumerate}
    \item Créer un motif de base (formes, couleurs).
    \item Sélectionner le motif.
    \item Utiliser \texttt{Objet > Motif > Créer}.
    \item La boîte de dialogue \emph{Options de motif} s'ouvre.
    \item Choisir le type de motif (côté, hexagone, etc.).
    \item Ajuster les options (taille, espacement, etc.).
    \item Cliquer sur \emph{Terminé}.
    \item Le motif est ajouté au nuancier.
\end{enumerate}

\subsubsection{Les options}

\begin{itemize}
    \item \terme{Type de motif} : côté, hexagone, etc.
    \item \terme{Taille de la tuile} : largeur et hauteur de la répétition.
    \item \terme{Espacement} : espace entre les tuiles.
    \item \terme{Décalage} : décalage des tuiles.
    \item \terme{Options de chevauchement} : chevauchement des tuiles.
\end{itemize}

\subsection{Appliquer un motif}

\begin{enumerate}
    \item Sélectionner l'objet.
    \item Ouvrir le panneau \emph{Nuancier}.
    \item Cliquer sur le motif (dans la section Motifs).
    \item Le motif est appliqué comme remplissage (ou contour).
\end{enumerate}

\subsection{Modifier un motif}

\begin{enumerate}
    \item Double-cliquer sur le motif dans le nuancier.
    \item La boîte de dialogue \emph{Options de motif} s'ouvre.
    \item Modifier les options.
    \item Cliquer sur \emph{Terminé}.
\end{enumerate}

\subsection{Les utilisations}

Les motifs sont utilisés pour :

\begin{itemize}
    \item \terme{Fonds} : créer des textures.
    \item \terme{Formes} : décorer des formes.
    \item \terme{Identités} : créer des éléments graphiques.
    \item \terme{Illustrations} : enrichir les compositions.
    \item \terme{Packagings} : créer des motifs de marque.
\end{itemize}

\subsection{Les bonnes pratiques}

\begin{itemize}
    \item \terme{Commencer} par un motif simple.
    \item \terme{Ajuster} la taille et l'espacement.
    \item \terme{Tester} le motif à différentes échelles.
    \item \terme{Enregistrer} les motifs dans le nuancier.
    \item \terme{Utiliser} les motifs pour créer des textures.
    \item \terme{Vérifier} le résultat.
\end{itemize}

\subsection*{Conclusion de la section}

Un \terme{motif} est une image répétée. On le crée avec \texttt{Objet > Motif > Créer}. On peut choisir le type, la taille, l'espacement, le décalage. On l'applique comme remplissage ou contour, via le nuancier. On peut le modifier. Les motifs sont utilisés pour créer des fonds, des textures, des décors. Il est essentiel de commencer simple, d'ajuster les options, de tester, de vérifier.

Maintenant que nous savons créer des motifs, nous allons, dans la section suivante, apprendre à \terme{gérer les transparences}.

% ============================================================
\section{Gérer les transparences}
% ============================================================

\subsection{Qu'est-ce que la transparence ?}

\begin{quote}
\textbf{La transparence est la propriété d'un objet qui laisse voir ce qui est derrière lui. Elle est définie par un pourcentage d'opacité : 100\% = opaque, 0\% = transparent.}
\end{quote}

La transparence est utilisée pour :

\begin{itemize}
    \item \terme{Superposer} des objets.
    \item \terme{Créer} des effets de profondeur.
    \item \terme{Adoucir} des transitions.
    \item \terme{Enrichir} les compositions.
    \item \terme{Créer} des ambiances.
\end{itemize}

\subsection{Le panneau Transparence}

\subsubsection{Accès au panneau}

Le panneau \terme{Transparence} (\texttt{Fenêtre > Transparence}) permet de gérer l'opacité et les modes de fusion.

\subsubsection{Les options}

\begin{itemize}
    \item \terme{Opacité} : pourcentage d'opacité (0\% à 100\%).
    \item \terme{Mode de fusion} : mode de fusion avec les objets du dessous.
    \item \terme{Masque d'opacité} : masque d'opacité.
    \item \terme{Isoler le groupe} : isoler le groupe pour les modes de fusion.
\end{itemize}

\subsection{Modifier l'opacité}

\begin{enumerate}
    \item Sélectionner l'objet.
    \item Ouvrir le panneau \emph{Transparence}.
    \item Saisir un pourcentage d'opacité.
    \item L'objet devient transparent.
\end{enumerate}

\subsection{Les modes de fusion}

\subsubsection{Définition}

Le \terme{mode de fusion} détermine comment les couleurs de l'objet se mélangent avec celles des objets du dessous. Les modes sont les mêmes que dans Photoshop : Normal, Multiplier, Écran, Superposition, etc.

\subsubsection{Utilisation}

\begin{enumerate}
    \item Sélectionner l'objet.
    \item Ouvrir le panneau \emph{Transparence}.
    \item Choisir un mode de fusion dans le menu déroulant.
    \item L'objet se mélange avec le fond.
\end{enumerate}

\subsection{Le masque d'opacité}

\subsubsection{Définition}

Le \terme{masque d'opacité} est un masque qui contrôle la transparence d'un objet en fonction de ses niveaux de gris (noir = transparent, blanc = opaque).

\subsubsection{Utilisation}

\begin{enumerate}
    \item Créer un objet (le masque).
    \item Sélectionner l'objet et le masque.
    \item Ouvrir le panneau \emph{Transparence}.
    \item Cliquer sur \emph{Créer un masque d'opacité}.
    \item Le masque est appliqué.
\end{enumerate}

\subsection{Les bonnes pratiques}

\begin{itemize}
    \item \terme{Utiliser} l'opacité pour superposer.
    \item \terme{Utiliser} les modes de fusion pour mélanger.
    \item \terme{Utiliser} les masques d'opacité pour les effets.
    \item \terme{Vérifier} le résultat.
    \item \terme{Attention} à l'impression : certaines transparences peuvent poser problème.
    \item \terme{Aplatir} les transparences si nécessaire pour l'impression.
\end{itemize}

\subsection*{Conclusion de la section}

La \terme{transparence} permet de superposer, de créer des effets de profondeur, d'adoucir, d'enrichir. On la gère avec le panneau Transparence : opacité, modes de fusion, masques d'opacité. Il est essentiel de l'utiliser avec parcimonie, de vérifier le résultat, de faire attention à l'impression.

Maintenant que nous savons gérer les transparences, nous allons, dans la dernière section de ce chapitre, apprendre à \terme{préparer les couleurs pour l'impression}.

% ============================================================
\section{Préparer les couleurs pour l'impression}
% ============================================================

\subsection{Pourquoi préparer les couleurs ?}

La \terme{préparation des couleurs} pour l'impression est essentielle pour :

\begin{itemize}
    \item \terme{Éviter} les mauvaises surprises (couleurs fausses).
    \item \terme{Respecter} les contraintes de l'imprimeur.
    \item \terme{Obtenir} un résultat conforme à l'écran.
    \item \terme{Optimiser} le coût d'impression.
    \item \terme{Assurer} la qualité.
\end{itemize}

\subsection{Les modes colorimétriques}

\subsubsection{RVB}

Le mode \terme{RVB} (Rouge, Vert, Bleu) est utilisé pour l'écran. Il ne convient pas à l'impression professionnelle.

\subsubsection{CMJN}

Le mode \terme{CMJN} (Cyan, Magenta, Jaune, Noir) est utilisé pour l'impression. Il faut convertir les documents en CMJN avant l'impression.

\subsubsection{Les tons directs}

Les \terme{tons directs} (Pantone) sont des encres spéciales, en plus des quatre encres CMJN. Ils sont utilisés pour les couleurs de marque.

\subsection{Convertir en CMJN}

\begin{enumerate}
    \item Ouvrir le document.
    \item Utiliser \texttt{Fichier > Document > Mode colorimétrique > Couleurs CMJN}.
    \item Les couleurs sont converties.
    \item Vérifier les couleurs (certaines peuvent changer).
\end{enumerate}

\subsection{Les profils colorimétriques}

\subsubsection{Définition}

Un \terme{profil colorimétrique} décrit la manière dont un appareil reproduit les couleurs. Il permet de convertir les couleurs entre appareils.

\subsubsection{Les profils pour l'impression}

\begin{itemize}
    \item \terme{ISO Coated v2} : pour l'impression offset sur papier couché.
    \item \terme{ISO Uncoated} : pour l'impression sur papier non couché.
    \item \terme{Coated FOGRA39} : pour l'impression offset en Europe.
    \item \terme{Profil de l'imprimeur} : à demander à l'imprimeur.
\end{itemize}

\subsection{Les fonds perdus}

\subsubsection{Définition}

Les \terme{fonds perdus} (ou \emph{bleed}) sont un dépassement de l'image au-delà des bords du document. Ils évitent les bords blancs après découpe.

\subsubsection{Utilisation}

\begin{enumerate}
    \item Ouvrir la boîte de dialogue \emph{Configuration du document}.
    \item Définir une valeur de fond perdu (par exemple, 3 mm).
    \item Étendre les éléments qui touchent les bords au-delà du document.
\end{enumerate}

\subsection{Les traits de coupe}

\subsubsection{Définition}

Les \terme{traits de coupe} (ou \emph{crop marks}) sont des repères qui indiquent à l'imprimeur où couper. Ils sont placés à l'extérieur des fonds perdus.

\subsubsection{Utilisation}

\begin{enumerate}
    \item Lors de l'export PDF, cocher \emph{Traits de coupe}.
    \item Les traits sont ajoutés au fichier.
\end{enumerate}

\subsection{L'export PDF pour l'impression}

\begin{enumerate}
    \item Utiliser \texttt{Fichier > Enregistrer sous} ou \texttt{Fichier > Exporter}.
    \item Choisir le format \emph{Adobe PDF}.
    \item Choisir le préréglage \emph{PDF/X-1a} ou \emph{PDF/X-4}.
    \item Vérifier les options (fonds perdus, traits de coupe, polices).
    \item Enregistrer.
\end{enumerate}

\subsection{Les bonnes pratiques}

\begin{itemize}
    \item \terme{Convertir} en CMJN avant l'impression.
    \item \terme{Utiliser} le profil de l'imprimeur.
    \item \terme{Prévoir} les fonds perdus.
    \item \terme{Ajouter} les traits de coupe.
    \item \terme{Vérifier} les couleurs (épreuve).
    \item \terme{Incorporer} les polices.
    \item \terme{Exporter} en PDF/X.
    \item \terme{Communiquer} avec l'imprimeur.
\end{itemize}

\subsection*{Conclusion de la section}

La \terme{préparation des couleurs} pour l'impression est essentielle. On convertit les documents en CMJN, on utilise le profil de l'imprimeur, on prévoit les fonds perdus, on ajoute les traits de coupe, on vérifie les couleurs (épreuve), on incorpore les polices, on exporte en PDF/X. Il est essentiel de communiquer avec l'imprimeur, de vérifier le résultat, de faire une épreuve.

% ============================================================
\section*{Conclusion du chapitre}
\addcontentsline{toc}{section}{Conclusion du chapitre}
% ============================================================

Nous voici au terme du chapitre 18. Nous avons appris à maîtriser les couleurs, les dégradés et les contours dans Illustrator :

\begin{itemize}
    \item \terme{Appliquer des couleurs} : remplissage et contour, panneaux Couleur et Nuancier, raccourcis.
    \item \terme{Nuanciers} : bibliothèques de couleurs, groupes, bibliothèques prédéfinies.
    \item \terme{Couleurs globales} : création, modification, tons directs.
    \item \terme{Dégradés linéaires} : création, modification, transparence.
    \item \terme{Dégradés radiaux} : création, centre, rayon.
    \item \terme{Contours} : épaisseur, pointillés, profil, flèches.
    \item \terme{Extrémités et angles} : bout, rond, carré ; angle, arrondi, biseau ; limite d'angle.
    \item \terme{Motifs} : création, application, modification.
    \item \terme{Transparences} : opacité, modes de fusion, masques d'opacité.
    \item \terme{Préparation pour l'impression} : CMJN, profils, fonds perdus, traits de coupe, PDF/X.
\end{itemize}

La couleur est au cœur du design graphique. Dans Illustrator, elle s'applique aux objets de deux manières : le remplissage et le contour. On peut utiliser des couleurs unies, des dégradés, des motifs, des transparences. On peut enregistrer des couleurs globales, créer des palettes cohérentes, préparer les couleurs pour l'impression. Maîtriser la couleur dans Illustrator, c'est maîtriser l'un des leviers les plus puissants du design graphique.

Mais la couleur ne s'utilise pas au hasard. Elle doit respecter les principes étudiés dans la première partie de ce cours (théorie des couleurs, harmonies, psychologie, contraste, lisibilité). Elle doit être adaptée au message, au public, au support. Elle doit être cohérente sur l'ensemble d'un projet. Elle doit être préparée correctement pour l'impression ou pour l'écran.

Ce chapitre t'a donné les bases des couleurs, des dégradés et des contours. Mais la maîtrise vient avec la pratique. N'hésite pas à expérimenter, à créer des palettes, à tester des dégradés, à préparer des fichiers pour l'impression. Chaque projet est une occasion de progresser.

Dans le chapitre suivant, nous aborderons la \terme{typographie et le lettrage dans Illustrator}. Tu apprendras à créer et modifier du texte, à comprendre les propriétés typographiques, à régler le crénage et l'interlignage, à créer des compositions typographiques, à utiliser le texte sur un tracé, à vectoriser du texte, à personnaliser des lettres, à créer un lettrage simple, à concevoir un titre publicitaire, et à préserver les versions modifiables des textes. Ces compétences te permettront de maîtriser la typographie dans Illustrator.

\vspace{2em}

\begin{center}
\textit{Fin du chapitre 18 --- Couleurs, dégradés et contours.}

\vspace{1em}

\textbf{Prochain chapitre : Chapitre 19 --- Typographie et lettrage dans Illustrator.}
\end{center}